% Generated by pipeline/gen_tex.py -- do not edit by hand.
\subsection{Ensemble of the Best-Performing Models}
\label{sec:ensemble}

Section~\ref{sec:results} showed that the three paradigms make qualitatively different errors, which is the precondition for a useful ensemble. We searched 1506 candidate combinations (hard and soft voting, uniform and validation-F1 weighting, 3--9 members) \textbf{on the validation split}, then scored the single winning configuration on the test split once. Selecting members by test performance would report a gain that does not generalise.

The selected system is soft voting with val-F1 weighting over 3 members: \textbf{BERT Base}, \textbf{Random Forest}, \textbf{Naive Bayes}.

\begin{table}[t]
\centering
\small
\begin{tabular}{lr}
\toprule
\textbf{System} & \textbf{Test Macro F1} \\
\midrule
BERT Base (best single) & 0.7647 \\
Ensemble & \textbf{0.7792} \\
\midrule
Difference & +1.45 pp \\
\bottomrule
\end{tabular}
\caption{Ensemble versus the best single model. Members were selected on validation (macro F1 0.7795) and the test split was scored once.}
\label{tab:ensemble}
\end{table}

\begin{table}[t]
\centering
\small
\begin{tabular}{lrr}
\toprule
\textbf{Member} & \textbf{Disagree} & \textbf{Own F1} \\
\midrule
Random Forest & 14.0\% & 0.6920 \\
Naive Bayes & 13.7\% & 0.7211 \\
\bottomrule
\end{tabular}
\caption{Ensemble member diversity: fraction of test documents on which each member disagrees with BERT Base. Decorrelation, not individual strength, is what earns a place in the vote.}
\label{tab:diversity}
\end{table}

\begin{table}[t]
\centering
\small
\begin{tabular}{lrrr}
\toprule
\textbf{Class} & \textbf{BERT Base} & \textbf{Ens.} & $\Delta$ \\
\midrule
Payday / title / personal loan & 0.5401 & 0.5735 & +3.34 \\
Vehicle / consumer loan & 0.6391 & 0.6554 & +1.63 \\
Money transfer / virtual currency & 0.7454 & 0.7608 & +1.54 \\
Student loan & 0.7912 & 0.8202 & +2.90 \\
Bank account or service & 0.8157 & 0.8292 & +1.34 \\
Mortgage & 0.8981 & 0.9004 & +0.23 \\
Credit card / prepaid card & 0.7679 & 0.7762 & +0.83 \\
Debt collection & 0.7693 & 0.7767 & +0.74 \\
Credit reporting & 0.9156 & 0.9204 & +0.48 \\
\bottomrule
\end{tabular}
\caption{Per-class F1, ordered from rarest to most common class.}
\label{tab:ensembleperclass}
\end{table}

The gain concentrates on the rare classes (+2.35 pp mean over the four smallest, against +0.57 pp over the four largest), which is the desirable direction: rare classes are where individual models are least confident and a vote has most to correct.

An ensemble pays the inference cost of every member: 268.3\,s over the test set against 267.3\,s for BERT Base alone (1.00$\times$).

That overhead is negligible, because the two added members are the cheapest models in the study: the ensemble costs 1.0\,s more than BERT Base alone across 303{,}213 documents, a 0.4\% increase, for +1.45 pp of Macro F1. Unlike the accuracy/latency trade of Table~\ref{tab:main}, there is no trade-off to weigh here: the ensemble dominates its strongest member at essentially equal cost.
